\chapter{设计的主体内容}

在着手进行上机设计之前首先做好大量准备：应熟悉课题，进行调查研究，收集国内、外资料并进行分析研究；同时还应完成交互界面的总体构思与实现方案设计。

\section{系统结构的设计}

\section{交互界面的设计和实现}

交互界面的设计应遵循简洁、直观、易用的原则，方便用户操作和观察算法执行过程。

\begin{equation}
  i(t)=I_{\mathrm{m}} e^{-\alpha t}\sin(\omega t+\varphi)
\label{equation:3-1}
\end{equation}

由公式（\ref{equation:3-1}）可以看出，电流随时间呈指数衰减的正弦变化规律。

\section{线性表的 OOP 设计}

计算机内部可以采用两种不同方法来表示一个线性表，它们分别是顺序表示法和链表表示法。

\begin{figure}[htbp]
  \centering
  \includegraphics[width=.72\linewidth]{img/example.pdf}
  \caption{过阻尼响应}
  \label{fig:3-1}
\end{figure}

过阻尼响应如图 \ref{fig:3-1} 所示，此时系统无振荡，输出缓慢趋近于稳态值。

\subsection{线性表的顺序存储的实现}

以上是顺序表的实现过程，第 1--16 行包含了 list 类的说明，接下来是成员函数的定义。

\subsection{线性表的链表存储的实现}

链表的实现包括两个类定义，第一个是 link 类，第二个是 list 类。由于一个链表由若干个单独的链结点对象组成，因此一个链结点应当作为单独的 link 类实现。

为使本模板在页眉页脚、图表编号和跨页排版方面能够与原始 Word 模板逐页对照，本章在后续页面中继续保留实现说明示例。实际撰写时，可继续补充系统流程、类图、测试步骤或核心代码设计说明，并根据篇幅需要增减内容。
